% =====================================================================
%  Informe - Grupo 3 - Desarrollo de Videojuegos (ESPE)
%  Unidad 2: Laberinto con IA en Go (backend) y Angular (frontend)
%  Compilar:  pdflatex InformeLaberintoGoAngular.tex  (dos veces, para el índice)
% =====================================================================
\documentclass[12pt,a4paper]{article}

\usepackage[utf8]{inputenc}
\usepackage[T1]{fontenc}
\usepackage[spanish,es-tabla,es-noshorthands]{babel}
\usepackage{lmodern}
\usepackage[margin=2.5cm]{geometry}
\usepackage{graphicx}
\usepackage{xcolor}
\usepackage{booktabs}
\usepackage{longtable}
\usepackage{array}
\usepackage{enumitem}
\usepackage{listings}
\usepackage{float}
\usepackage{caption}
\usepackage{subcaption}
\usepackage{fancyhdr}
\usepackage{amssymb}
\usepackage{tikz}
\usetikzlibrary{shapes.geometric,shapes.multipart,arrows.meta,positioning,calc,babel}
\usepackage[hidelinks]{hyperref}

\graphicspath{{Img/}}

% ---------- colores y utilidades ----------
\definecolor{espe}{RGB}{0,102,51}
\definecolor{xcol}{RGB}{8,145,178}
\definecolor{ocol}{RGB}{219,39,119}
\definecolor{codebg}{RGB}{245,245,250}

% Captura con borde fino; el argumento opcional recibe las opciones de la imagen.
\newcommand{\captura}[2][width=0.99\linewidth]{%
 \begingroup\setlength{\fboxsep}{0pt}\setlength{\fboxrule}{0.4pt}%
 \fcolorbox{gray!50}{white}{\includegraphics[#1]{#2}}\endgroup}

\captionsetup{font=small,labelfont=bf,labelsep=period}
\captionsetup[sub]{font=footnotesize,labelfont=bf}

% Salida real de consola, guardada en Img/consola-*.txt.
\definecolor{consolabg}{RGB}{24,26,38}
\definecolor{consolafg}{RGB}{226,232,240}
\lstdefinestyle{consola}{basicstyle=\ttfamily\scriptsize\color{consolafg},
  backgroundcolor=\color{consolabg}, frame=single, rulecolor=\color{consolabg},
  framesep=4pt, xleftmargin=4pt, xrightmargin=4pt, aboveskip=0pt, belowskip=0pt,
  breaklines=true, breakindent=1.2em, columns=fullflexible, keepspaces=true}

\lstdefinelanguage{Go}{
  morekeywords={package,import,func,return,if,else,for,range,switch,case,default,type,struct,var,const,interface,map,go,defer,true,false,nil},
  sensitive=true, morecomment=[l]{//}, morestring=[b]"}
\lstset{basicstyle=\ttfamily\footnotesize, backgroundcolor=\color{codebg},
  keywordstyle=\color{xcol}\bfseries, commentstyle=\color{gray}\itshape,
  stringstyle=\color{ocol}, frame=single, rulecolor=\color{gray!40},
  breaklines=true, columns=fullflexible, showstringspaces=false,
  literate={á}{{\'a}}1 {é}{{\'e}}1 {í}{{\'i}}1 {ó}{{\'o}}1 {ú}{{\'u}}1 {ñ}{{\~n}}1
    {¡}{{!`}}1 {«}{{\guillemotleft}}1 {»}{{\guillemotright}}1
    {✔}{{$\checkmark$}}1 {▶}{{$\blacktriangleright$}}1 {ℹ}{{\textit{i}}}1}

\pagestyle{fancy}
\fancyhf{}
\fancyhead[L]{\small Desarrollo de Videojuegos -- Unidad 2}
\fancyhead[R]{\small Grupo N.º 3}
\fancyfoot[C]{\thepage}
\renewcommand{\headrulewidth}{0.4pt}
\setlength{\headheight}{15pt}

\begin{document}
\sloppy

% =====================================================================
%  CARÁTULA
% =====================================================================
\begin{titlepage}
  \centering
  \includegraphics[width=0.55\textwidth]{logoINFESPE.png}\par\vspace{1cm}
  {\Large\bfseries UNIVERSIDAD DE LAS FUERZAS ARMADAS -- ESPE\par}
  \vspace{0.3cm}
  {\large DEPARTAMENTO DE CIENCIAS DE LA COMPUTACIÓN\par}
  \vspace{0.2cm}
  {\large CARRERA DE INGENIERÍA EN SOFTWARE\par}
  \vspace{1.2cm}
  {\large\bfseries ASIGNATURA: DESARROLLO DE VIDEOJUEGOS\par}
  \vspace{0.3cm}
  {\large UNIDAD 2\par}
  \vspace{1.2cm}
  {\color{espe}\rule{\linewidth}{1pt}}\par\vspace{0.3cm}
  {\LARGE\bfseries Laberinto con Inteligencia Artificial\par}
  \vspace{0.2cm}
  {\large Modos manual, mano derecha y mano izquierda -- Go (Web) + Angular\par}
  \vspace{0.3cm}
  {\color{espe}\rule{\linewidth}{1pt}}\par
  \vspace{1.2cm}
  {\large DOCENTE: ING. CÉSAR VILLACÍS SILVA, MSC.\par}
  \vspace{0.8cm}
  {\large\bfseries ESTUDIANTES:\par}
  \vspace{0.2cm}
  {\large Diana Guerra\par Paul Moreno\par Germán Cáceres\par}
  \vspace{0.8cm}
  {\large\bfseries GRUPO N.º 3\par}
  \vfill
  {\large FECHA: NOVIEMBRE 2026\par}
\end{titlepage}

\tableofcontents
\newpage

% =====================================================================
\section{Presentación del proyecto}
% =====================================================================
Se migró a la web el videojuego \emph{Laberinto con Inteligencia Artificial} del
docente, desarrollado originalmente en C\# con Windows Forms. El jugador recorre
un laberinto de 19$\times$19 casillas desde el inicio hasta la meta, controlándolo
con el teclado o con la voz, o deja que la computadora lo resuelva siguiendo la
pared derecha o la izquierda.

La arquitectura es cliente--servidor: un backend en Go ejecuta toda la lógica del
juego y la expone en una API HTTP/JSON, y un frontend en Angular presenta el
tablero y envía las acciones del usuario. Existe una segunda migración del mismo
juego en JavaScript puro, documentada en un informe aparte.

\begin{table}[H]
\centering
\caption{Modos de juego implementados}
\begin{tabular}{p{2.8cm}p{5cm}p{5.6cm}}
\toprule
\textbf{Modo} & \textbf{Comportamiento} & \textbf{Formulario original}\\
\midrule
Manual & El usuario mueve al caminante con las flechas o con la voz. & \texttt{frmBeginnerManual}\\
Mano derecha & La IA sigue la pared derecha hasta la meta. & \texttt{frmBeginnerAutomatic}\\
Mano izquierda & La IA sigue la pared izquierda hasta la meta. & \texttt{frmBeginnerAutomaticLeft}\\
\bottomrule
\end{tabular}
\end{table}

Se conservaron las dos clases del original como estructuras de Go,
\texttt{Walker} (11 métodos) y \texttt{NonPlayerCharacter} (19 métodos), con los
mismos nombres, condiciones y orden de evaluación. La lógica no se reemplazó por
tablas, rutas precalculadas ni otro algoritmo.

% =====================================================================
\section{Especificación de Requerimientos (IEEE 830)}
% =====================================================================

\subsection{Introducción}

\subsubsection{Propósito}
Se especifican los requisitos funcionales y no funcionales del videojuego web
\emph{Laberinto con IA} en Go y Angular, con base en IEEE 830-1998~\cite{ieee830}.

\subsubsection{Alcance}
El sistema presenta un menú con tres modos y un laberinto de nivel principiante.
En el modo manual el usuario controla al caminante; en los modos automáticos un
\emph{Personaje No Jugador} (NPC) aplica la regla de la mano derecha o de la mano
izquierda. El backend guarda una partida por pestaña del navegador y responde a
cada acción con el estado del tablero.

\subsubsection{Definiciones, acrónimos y abreviaturas}
\begin{description}[leftmargin=2.8cm,style=nextline]
  \item[Walker] Estructura del caminante controlado por el usuario.
  \item[NPC] \emph{Non-Player Character}: estructura \texttt{NonPlayerCharacter}, que resuelve el laberinto sola.
  \item[Matriz] Tablero de 19$\times$19 con 0 (pared), 1 (camino), 2 (inicio) y 3 (meta).
  \item[Tick] Cada avance de la IA, que el frontend solicita 50 ms después de recibir la respuesta anterior.
  \item[Sesión] Identificador aleatorio de cada pestaña, con el que el backend separa las partidas.
\end{description}

\subsubsection{Referencias}
\begin{itemize}
  \item IEEE Std 830-1998, \emph{Recommended Practice for Software Requirements Specifications}.
  \item Proyecto C\# Windows Forms del docente: clases \texttt{Walker} y \texttt{NonPlayerCharacter}, y formularios del nivel principiante.
  \item Referencias del C\# original usadas en la migración: \texttt{testdata/original\_decisions.json} y \texttt{original\_routes.json}.
\end{itemize}

\subsection{Descripción general}

\subsubsection{Perspectiva del producto}
El backend (\texttt{backend/}) contiene la matriz, \texttt{Walker},
\texttt{NonPlayerCharacter} y el servidor HTTP, solo con la librería estándar de
Go. El frontend (\texttt{frontend/src/main.ts}) es un componente Angular que
dibuja el tablero con la matriz recibida, ubica al caminante según su posición y
programa los ticks. Durante el desarrollo, Angular redirige \texttt{/api} al
puerto 8080 con \texttt{proxy.conf.json}.

\subsubsection{Funciones del producto}
Se permite elegir el modo, mover al caminante con las flechas o con las órdenes de
voz \emph{right}, \emph{left}, \emph{up} y \emph{down}, resolver el laberinto con
la IA, reiniciar el recorrido, volver al menú y salir. Al llegar a la meta se
muestra el mensaje que devuelve el backend.

\subsubsection{Características de los usuarios}
Estudiantes y jugadores sin conocimiento técnico. El modo manual requiere usar
las flechas del teclado; la voz es opcional.

\subsubsection{Restricciones}
\begin{itemize}
  \item La migración debe conservar los métodos, condiciones y orden de evaluación del C\#.
  \item Backend en Go~1.22 o superior, solo con la librería estándar.
  \item Frontend en Angular~19 con Node.js~22 o~24.
  \item La voz depende de la API \texttt{SpeechRecognition} del navegador y del permiso de micrófono.
\end{itemize}

\subsubsection{Suposiciones y dependencias}
El backend escucha en \texttt{127.0.0.1:8080} y el frontend en
\texttt{http://localhost:4200}; ambos deben estar en ejecución.

\subsection{Requerimientos específicos}

\subsubsection{Requerimientos funcionales}
\begin{longtable}{p{1.3cm}p{9.2cm}p{2.6cm}}
\caption{Requerimientos funcionales}\\
\toprule
\textbf{ID} & \textbf{Descripción} & \textbf{Prioridad}\\
\midrule\endfirsthead
\toprule
\textbf{ID} & \textbf{Descripción} & \textbf{Prioridad}\\
\midrule\endhead
RF01 & El sistema usará la misma matriz de 19$\times$19 del original, con inicio en (1,2) y meta en (17,16). & Alta\\
RF02 & El menú permitirá elegir el modo manual, mano derecha o mano izquierda, y salir del juego. & Alta\\
RF03 & En modo manual, las flechas moverán al caminante una casilla; una pared impedirá el movimiento. & Alta\\
RF04 & En modo manual, las órdenes de voz moverán al caminante; otros textos se añadirán al cuadro de reconocimiento. & Media\\
RF05 & Al llegar a la meta con el teclado se mostrará «Felicitaciones, buen trabajo.» y el teclado dejará de mover al caminante. & Alta\\
RF06 & La acción \texttt{solve} reiniciará la posición, determinará la dirección inicial y hará el primer movimiento de la IA. & Alta\\
RF07 & Cada acción \texttt{tick} ejecutará un movimiento de la IA y comprobará la meta. & Alta\\
RF08 & La IA de mano derecha y la de mano izquierda conservarán las condiciones del original, incluidos los \texttt{if} independientes de la mano derecha. & Alta\\
RF09 & Al llegar a la meta en modo automático, el backend devolverá \texttt{running = false} y el mensaje, y el frontend dejará de pedir ticks. & Alta\\
RF10 & La acción \texttt{exit} liberará la sesión; un \texttt{tick} posterior responderá con un error. & Media\\
\bottomrule
\end{longtable}

\subsubsection{Requerimientos no funcionales}
\begin{longtable}{p{1.3cm}p{2.9cm}p{9.2cm}}
\caption{Requerimientos no funcionales}\\
\toprule
\textbf{ID} & \textbf{Tipo} & \textbf{Descripción}\\
\midrule\endfirsthead
\toprule
\textbf{ID} & \textbf{Tipo} & \textbf{Descripción}\\
\midrule\endhead
RNF01 & Fidelidad & Mismos métodos de \texttt{Walker} y \texttt{NonPlayerCharacter}, sin métodos de juego añadidos.\\
RNF02 & Rendimiento & Cada petición de tick se responderá muy por debajo del intervalo de 50 ms.\\
RNF03 & Robustez & La API rechazará métodos, acciones, modos y sesiones inválidos, y limitará el cuerpo de la petición a 4\,096 bytes.\\
RNF04 & Concurrencia & Varias pestañas podrán jugar a la vez, cada una con su sesión, protegidas por exclusión mutua.\\
RNF05 & Portabilidad & La interfaz se adaptará a celular sin desplazamiento horizontal.\\
RNF06 & Calidad & Las pruebas de paridad deberán aprobar con \texttt{go test ./...} y \texttt{go vet ./...} no emitirá diagnósticos.\\
\bottomrule
\end{longtable}

% =====================================================================
\section{Modelo de Inteligencia Artificial}
% =====================================================================
El laberinto se representa con una matriz de 19$\times$19. Cada paso mueve al
caminante una casilla y desplaza su posición 25 píxeles lógicos.

\begin{table}[H]
\centering
\caption{Valores de la matriz del laberinto}
\begin{tabular}{cll}
\toprule
\textbf{Valor} & \textbf{Significado} & \textbf{Casillas}\\
\midrule
0 & Pared & 232\\
1 & Camino & 127\\
2 & Inicio, en (1,2) & 1\\
3 & Meta, en (17,16) & 1\\
\bottomrule
\end{tabular}
\end{table}

La IA aplica la regla de \emph{seguir la pared}: si se mantiene siempre una pared
del mismo lado, el caminante termina recorriendo el borde de la región hasta la
salida. \texttt{NonPlayerCharacter} observa las cuatro casillas vecinas con
\texttt{ObserveRightSide}, \texttt{ObserveLeftSide}, \texttt{ObserveUpSide} y
\texttt{ObserveDownSide}, y decide según la dirección actual. Por ejemplo, la mano
derecha mirando a la derecha, en Go:

\begin{lstlisting}[language=Go]
if p.Direction == Right {
    if !p.ObserveUpSide() && !p.ObserveDownSide() && p.ObserveRightSide() {
        p.Direction = Right
        p.MoveRightSide()
    } else if p.ObserveUpSide() && !p.ObserveDownSide() && p.ObserveRightSide() {
        p.Direction = Right
        p.MoveRightSide()
    } else if p.ObserveDownSide() {
        p.Direction = Down
        p.MoveDownSide()
    } // ... gira hacia arriba o retrocede si no hay salida
}
if p.Direction == Left { // if independiente, como en el C#
\end{lstlisting}

En la mano derecha, los bloques de cada dirección son \texttt{if} independientes,
como en el C\#: si un bloque cambia la dirección, el siguiente puede volver a mover
al caminante en la misma llamada. En la mano izquierda los bloques son
\texttt{else if} y hay un solo movimiento por llamada. Se conservó esa diferencia.

\begin{table}[H]
\centering\small
\caption{Recorridos de la IA en el laberinto principiante}
\begin{tabular}{lrrrrr}
\toprule
\textbf{Modo} & \textbf{Ticks} & \textbf{Movim.} & \shortstack{\textbf{Ticks con}\\\textbf{2+ movim.}} & \shortstack{\textbf{Casillas}\\\textbf{distintas}} & \shortstack{\textbf{Tiempo}\\\textbf{nominal}}\\
\midrule
Mano izquierda & 66 & 68 & 1 & 60 & 3,35 s\\
Mano derecha & 161 & 178 & 15 & 106 & 8,10 s\\
\bottomrule
\end{tabular}
\end{table}
El tiempo nominal es de 50 ms por tick; en esta versión se suma la latencia de
cada petición, porque el siguiente tick se programa al recibir la respuesta. El
camino más corto del inicio a la meta tiene 54 movimientos; ninguna de las dos
manos lo sigue, solo garantizan llegar a la salida en un laberinto como este.

% =====================================================================
\section{Diagramas}
% =====================================================================
Las figuras~\ref{fig:uml-casos}, \ref{fig:uml-secuencia} y \ref{fig:uml-clases}
representan los casos de uso, la resolución automática a través de la API y la
estructura del sistema.

\subsection{Diagrama de casos de uso}
\begin{figure}[H]
\centering
\begin{tikzpicture}[
 uc/.style={ellipse,draw,fill=xcol!8,minimum width=4.4cm,minimum height=0.9cm,font=\small,align=center},
 sysuc/.style={uc,fill=ocol!8,minimum width=3.8cm},
 rel/.style={dashed,->},
 lbl/.style={font=\scriptsize,fill=white,inner sep=1pt},>=Stealth]
 \node[draw,rounded corners,minimum width=11.5cm,minimum height=10cm] (sys) at (1.15,0) {};
 \node[font=\bfseries\small,anchor=south] at (sys.north) {Laberinto con IA (Go + Angular)};
 \node[uc] (mode) at (-2,3.7) {Elegir modo};
 \node[uc] (keys) at (-2,1.9) {Mover con el teclado};
 \node[uc] (voice) at (-2,0) {Mover con la voz};
 \node[uc] (solve) at (-2,-1.9) {Resolver con la IA};
 \node[uc] (exit) at (-2,-3.7) {Volver al menú o salir};
 \node[sysuc] (move) at (4.7,2.6) {Mover caminante};
 \node[sysuc] (goal) at (4.7,0) {Comprobar meta};
 \node[sysuc] (ai) at (4.7,-2.6) {Seguir la pared};
 \begin{scope}[shift={(-6.7,0)},line width=0.9pt,xcol!70!black]
  \draw[fill=xcol!8] (0,0.8) circle (0.27);
  \draw (0,0.53) -- (0,-0.3) (-0.42,0.2) -- (0.42,0.2) (0,-0.3) -- (-0.33,-0.9) (0,-0.3) -- (0.33,-0.9);
 \end{scope}
 \node[font=\small] at (-6.7,-1.2) {Usuario};
 \coordinate (usr) at (-6.1,0);
 \foreach \n in {mode,keys,voice,solve,exit} \draw (usr) -- (\n.west);
 \draw[rel] (keys.east) -- node[lbl]{$\ll$include$\gg$} (move.west);
 \draw[rel] (voice.north east) -- node[lbl]{$\ll$include$\gg$} (move.south west);
 \draw[rel] (move.south) -- node[lbl]{$\ll$include$\gg$} (goal.north);
 \draw[rel] (solve.east) -- node[lbl]{$\ll$include$\gg$} (ai.west);
 \draw[rel] (ai.north) -- node[lbl]{$\ll$include$\gg$} (goal.south);
\end{tikzpicture}
\caption{Casos de uso. El NPC y el backend forman parte del sistema: el usuario solo inicia la resolución.}
\label{fig:uml-casos}
\end{figure}

\subsection{Diagrama de secuencia: resolver con la mano derecha}
\begin{figure}[H]
\centering\resizebox{\textwidth}{!}{%
\begin{tikzpicture}[font=\small,>=Stealth,
 obj/.style={draw,fill=xcol!10,minimum width=2.8cm,minimum height=1.2cm,align=center},
 msg/.style={align=center,fill=white,inner sep=1pt},
 self/.style={right,align=left}]
 \node[obj] (a) at (0,0) {:Usuario};
 \node[obj] (b) at (4.2,0) {:App\\(Angular)};
 \node[obj] (c) at (8.6,0) {:API Go\\POST /api/game};
 \node[obj] (d) at (15.6,0) {:NonPlayer-\\Character};
 \foreach \n in {a,b,c,d} \draw[dashed,gray] (\n.south) -- ++(0,-12.8);
 \foreach \x/\ini/\fin in {4.2/-1.5/-12.4,8.6/-2.3/-5.6,15.6/-3.1/-4.8,8.6/-7.6/-10.6,15.6/-8.4/-9.8}
  \draw[fill=xcol!12,draw=xcol!60!black] (\x-0.09,\ini) rectangle (\x+0.09,\fin);
 \draw[->] (0,-1.5) -- node[msg,above=1pt]{clic en Resolver} (4.2,-1.5);
 \draw[->] (4.2,-2.3) -- node[msg,above=1pt]{action: solve} (8.6,-2.3);
 \draw[->] (8.6,-3.1) -- node[msg,above=1pt]{InitializeRoad, LocalizeStart} (15.6,-3.1);
 \draw[->] (8.6,-4.0) -- node[msg,above=1pt]{DetermineStartDirection,\\RapidMovementWithRightHandAI} (15.6,-4.0);
 \draw[dashed,->] (8.6,-5.4) -- node[msg,above=1pt]{JSON: posición, running} (4.2,-5.4);
 \draw[gray,rounded corners=1pt] (3.5,-6.4) rectangle (17.0,-11.2);
 \node[anchor=north west,draw=gray,fill=white,font=\scriptsize,inner sep=2.5pt] at (3.5,-6.4) {\textbf{loop} [50 ms después de cada respuesta, mientras running]};
 \draw[->] (4.2,-7.6) -- node[msg,above=1pt]{action: tick} (8.6,-7.6);
 \draw[->] (8.6,-8.4) -- node[msg,above=1pt]{RapidMovementWithRightHandAI} (15.6,-8.4);
 \draw[->] (8.6,-9.2) -- node[msg,above=1pt]{DetermineFinishDirection} (15.6,-9.2);
 \draw[dashed,->] (8.6,-10.4) -- node[msg,above=1pt]{JSON: posición, running, message} (4.2,-10.4);
 \draw[dashed,->] (4.2,-12.2) -- node[msg,above=1pt]{tablero y aviso de meta} (0,-12.2);
\end{tikzpicture}}
\caption{Resolución automática con la mano derecha. El backend ejecuta la lógica;
Angular solo programa los ticks y dibuja el estado. La mano izquierda sigue el
mismo flujo con \texttt{RapidMovementWithLeftHandAI}.}
\label{fig:uml-secuencia}
\end{figure}

\subsection{Diagrama de clases}
\begin{figure}[H]
\centering\resizebox{\textwidth}{!}{%
\begin{tikzpicture}[font=\scriptsize,>=Stealth,
 cls/.style={draw,rectangle split,rectangle split parts=3,rectangle split part fill={xcol!14,white,white},text width=6.6cm,align=left,inner sep=3pt,anchor=north}]
 \node[cls,text width=6cm] (walker) at (-7.3,8) {\makebox[\linewidth]{\textbf{Walker}}
  \nodepart{second}+Matrix: [][]int\\+Row, Column, X, Y: int\\+Table: []string\\+Data, Message: string
  \nodepart{third}+CreateMatrix(r, c), InitializeMatrix(a)\\+InitializeRoad(), LocalizeStart()\\+MoveRight/Left/Up/DownSide()\\+InitializeGrammar()\\+ControlMovesInLaberynth(key string)\\+ControlSpeechMoves(text string)};
 \node[cls] (npc) at (0,8) {\makebox[\linewidth]{\textbf{NonPlayerCharacter}}
  \nodepart{second}+Matrix: [][]int\\+Row, Column, X, Y: int\\+Direction, DirectionText, Message: string\\+Miliseconds: int
  \nodepart{third}+CreateMatrix, InitializeMatrix\\+InitializeRoad, LocalizeStart\\+MoveRight/Left/Up/DownSide()\\+DetermineStartDirection()\\+DetermineFinishDirection(): bool\\+ObserveRight/Left/Up/DownSide(): bool\\+RapidMovementWithLeftHandAI()\\+RapidMovementWithRightHandAI()\\+InitializeRapidMovement()\\+InitializePausedMovement()\\+PausedMovementWithAI()};
 \node[cls,rectangle split part fill={gray!15,white,white},text width=4.6cm] (game) at (-7.3,1.6) {\makebox[\linewidth]{\textbf{game}}
  \nodepart{second}+Mode: string\\+NPC: NonPlayerCharacter\\+Walker: Walker\\+Running: bool
  \nodepart{third}---};
 \node[cls,rectangle split part fill={gray!15,white,white},text width=5.4cm] (main) at (0,0.6) {\makebox[\linewidth]{$\ll$función$\gg$ \textbf{main} (handler /api/game)}
  \nodepart{second}games: map[string]*game\\mutex: sync.Mutex
  \nodepart{third}load, solve, tick, key, speech, exit};
 \node[cls,rectangle split part fill={ocol!12,white,white},text width=4.8cm] (app) at (7.3,8) {\makebox[\linewidth]{\textbf{App} (Angular)}
  \nodepart{second}mode, state, busy, session
  \nodepart{third}load(mode)\\solve()\\schedule()\\keydown(event)\\StartSpeech()\\exit()};
 \draw[{Diamond[fill=black]}-] (game.north) -- node[right]{1} (walker.south);
 \draw[{Diamond[fill=black]}-] (game.north east) -- node[above,sloped]{1} (npc.south west);
 \draw[{Diamond[open]}-] (main.west) -- node[above]{0..*} (game.east);
 \draw[dashed,->] (app.south) |- node[pos=0.25,right]{HTTP / JSON} (main.east);
\end{tikzpicture}}
\caption{Estructura en UML. \texttt{Walker} y \texttt{NonPlayerCharacter} son la
lógica traducida desde C\#; \texttt{game} guarda una partida por sesión y el
manejador de \texttt{main} despacha las acciones. \texttt{App} es el componente
Angular, que se comunica con el backend por HTTP.}
\label{fig:uml-clases}
\end{figure}

% =====================================================================
\section{Metodología SCRUM + XP}
% =====================================================================
El desarrollo se organizó en sprints, con Scrum como marco de planificación y
prácticas XP para el diseño, la implementación y las pruebas.

\subsection{Roles}
\begin{table}[H]
\centering
\caption{Distribución de roles y responsabilidades}
\begin{tabular}{ll}
\toprule
\textbf{Rol} & \textbf{Responsable}\\
\midrule
Product Owner & Ing. César Villacís Silva, MSc.\\
Scrum Master & Diana Guerra\\
Desarrollo, pruebas y documentación & Diana Guerra, Paul Moreno y Germán Cáceres\\
\bottomrule
\end{tabular}
\end{table}

\subsection{Product Backlog}
\begin{longtable}{p{1.1cm}p{8.6cm}p{1.2cm}p{1.6cm}}
\caption{Historias de usuario}\\
\toprule
\textbf{HU} & \textbf{Historia} & \textbf{Pts} & \textbf{Sprint}\\
\midrule\endfirsthead
\toprule
\textbf{HU} & \textbf{Historia} & \textbf{Pts} & \textbf{Sprint}\\
\midrule\endhead
HU01 & Como usuario, quiero ver el laberinto con el inicio y la meta, para saber adónde llegar. & 3 & 1\\
HU02 & Como usuario, quiero mover al caminante con las flechas, para recorrer el laberinto. & 5 & 1\\
HU03 & Como usuario, quiero que las paredes me detengan, para que el recorrido respete el laberinto. & 3 & 1\\
HU04 & Como usuario, quiero mover al caminante con mi voz, para jugar sin teclado. & 5 & 2\\
HU05 & Como usuario, quiero que la IA resuelva el laberinto con la mano derecha, para observar su estrategia. & 8 & 2\\
HU06 & Como usuario, quiero que la IA lo resuelva con la mano izquierda, para comparar ambas rutas. & 5 & 2\\
HU07 & Como usuario, quiero un aviso al llegar a la meta, para saber que terminé. & 2 & 3\\
HU08 & Como usuario, quiero que varias pestañas jueguen a la vez sin mezclarse, para compartir el equipo. & 3 & 3\\
\bottomrule
\end{longtable}

\subsection{Sprints}
\begin{description}
  \item[Sprint 1 -- Núcleo en Go.] Traducción de la matriz y de \texttt{Walker}; movimiento con teclado y paredes; pruebas de paridad.
  \item[Sprint 2 -- IA y API.] Traducción de \texttt{NonPlayerCharacter}, acciones \texttt{solve} y \texttt{tick}, sesiones por pestaña y voz.
  \item[Sprint 3 -- Angular e integración.] Componente Angular con interfaz arcade, proxy de desarrollo y validación en el navegador.
\end{description}

\subsection{Prácticas XP aplicadas}
\begin{itemize}
 \item \textbf{Diseño simple:} un solo endpoint con seis acciones y un solo componente Angular.
 \item \textbf{Pruebas automatizadas:} comparación de cada decisión y de cada recorrido con las referencias del C\#.
 \item \textbf{Estándares de código:} \texttt{go vet} sin diagnósticos y compilación de producción de Angular.
\end{itemize}

\subsection{Diseño}
\begin{table}[H]
\centering\small
\caption{Estructura del proyecto}
\begin{tabular}{p{5.4cm}p{8.2cm}}
\toprule
\textbf{Archivo} & \textbf{Responsabilidad}\\
\midrule
\texttt{backend/matrix.go} & Matriz original de 19$\times$19\\
\texttt{backend/walker.go} & \texttt{Walker}: modo manual con teclado y voz\\
\texttt{backend/npc.go} & \texttt{NonPlayerCharacter}: IA de mano derecha e izquierda\\
\texttt{backend/main.go} & Servidor HTTP, validación y sesiones\\
\texttt{backend/parity\_test.go} & Pruebas de paridad con el C\#\\
\texttt{frontend/src/main.ts} & Componente Angular: menú, tablero, ticks y voz\\
\texttt{frontend/src/styles.css} & Interfaz arcade adaptable\\
\texttt{frontend/proxy.conf.json} & Redirección de \texttt{/api} al backend\\
\bottomrule
\end{tabular}
\end{table}

\begin{table}[H]
\centering\small
\caption{Acciones de \texttt{POST /api/game}}
\begin{tabular}{p{1.8cm}p{5.8cm}p{5.6cm}}
\toprule
\textbf{Acción} & \textbf{Efecto} & \textbf{Equivale a (original)}\\
\midrule
load & Crea la partida del modo indicado & Abrir el formulario\\
solve & Reinicia y hace el primer movimiento de la IA & \texttt{btnSolve\_Click}\\
tick & Mueve la IA y comprueba la meta & \texttt{tmrBeginner\_Tick}\\
key & Mueve con una flecha & \texttt{frmLaberynth\_KeyDown}\\
speech & Mueve con una orden de voz & \texttt{SpeechRecognized}\\
exit & Libera la sesión & \texttt{btnExit\_Click}\\
\bottomrule
\end{tabular}
\end{table}

\subsection{Desarrollo}
\begin{table}[H]
\centering
\caption{Correspondencia C\# $\rightarrow$ Go y Angular}
\small
\begin{tabular}{ll}
\toprule
\textbf{C\# (original)} & \textbf{Go y Angular}\\
\midrule
\texttt{PictureBox picWalker} & campos \texttt{X} e \texttt{Y} de la estructura\\
\texttt{TextBox txtDirection / txtData} & campos \texttt{DirectionText} y \texttt{Data}\\
\texttt{enum Direction} & constantes \texttt{Up}, \texttt{Down}, \texttt{Left} y \texttt{Right}\\
\texttt{Timer tmrBeginner} (50 ms) & \texttt{setTimeout} en Angular y acción \texttt{tick}\\
\texttt{Thread.Sleep(mMiliseconds)} & \texttt{time.Sleep}\\
\texttt{System.Speech} & \texttt{SpeechRecognition} del navegador y acción \texttt{speech}\\
\texttt{MessageBox.Show} & campo \texttt{Message} en el JSON y aviso en la página\\
Formularios Windows Forms & componente Angular y manejador \texttt{/api/game}\\
\bottomrule
\end{tabular}
\end{table}

\subsection{Pruebas}
Se ejecutaron el 24 de noviembre de 2026, desde
\texttt{UNIDAD II/Projectos/Laberinto Go Angular/backend}, con
\texttt{go vet ./...} y \texttt{go test -count=1 -v ./...}. Las referencias son
archivos JSON generados con el C\# original durante la migración.

\begin{longtable}{p{1cm}p{3.6cm}p{3.3cm}p{3.2cm}p{1.6cm}}
\caption{Objetivos, escenarios y resultados de las pruebas automatizadas}\\
\toprule ID & Objetivo y procedimiento & Esperado & Obtenido & Estado\\\midrule\endfirsthead
\toprule ID & Objetivo y procedimiento & Esperado & Obtenido & Estado\\\midrule\endhead
G01 & Ejecutar 320 decisiones de referencia (64 por cada uno de 5 métodos). & Misma fila, columna, dirección, posición y resultado. & 320 coincidencias. & PASS\\
G02 & Reproducir los recorridos completos de ambas manos, tick por tick. & Cada posición igual a la del original. & Izquierda 66 ticks; derecha 161; meta (17,16). & PASS\\
G03 & Teclado, voz, paredes y llegada a la meta. & Movimientos válidos, bloqueo y mensaje. & Teclado bloqueado en la meta; voz no, como en el original. & PASS\\
G04 & Análisis estático con \texttt{go vet}. & Sin diagnósticos. & Sin diagnósticos. & PASS\\
\bottomrule
\end{longtable}

\begin{figure}[H]
\lstinputlisting[style=consola]{Img/consola-pruebas.txt}
\caption{Salida de consola de las pruebas de Go. G01--G03 corresponden a \texttt{TestOriginalDecisions}, \texttt{TestOriginalRoutes} y \texttt{TestWalkerControls}.}
\label{fig:consola-pruebas}
\end{figure}

\subsubsection{Pruebas de la API}
Con el backend y el frontend en ejecución, se enviaron peticiones con Node.js a
\texttt{http://localhost:4200/api/game}, a través del proxy de Angular.
\begin{longtable}{p{1cm}p{3.6cm}p{3.3cm}p{3.2cm}p{1.6cm}}
\caption{Pruebas de la API}\\
\toprule ID & Objetivo y procedimiento & Esperado & Obtenido & Estado\\\midrule\endfirsthead
\toprule ID & Objetivo y procedimiento & Esperado & Obtenido & Estado\\\midrule\endhead
A01 & Resolver con cada mano pidiendo ticks hasta que \texttt{running} sea falso. & 66 y 161 ticks; meta (17,16). & Iguales al original; dirección vacía. & PASS\\
A02 & Pedir un tick después de \texttt{exit}. & Sesión liberada. & HTTP 409. & PASS\\
A03 & Modo manual: flecha contra una pared y luego dos flechas válidas. & Sin movimiento y luego (2,2). & (1,2) y luego (2,2). & PASS\\
A04 & Modo, acción o sesión inválidos. & Rechazo. & HTTP 400 en los tres casos. & PASS\\
A05 & Petición con GET. & Método no permitido. & HTTP 405. & PASS\\
A06 & Medir la latencia de 227 ticks. & Muy por debajo de 50 ms. & Mediana 3,45 ms; P95 4,94 ms; máximo 7,40 ms. & PASS\\
\bottomrule
\end{longtable}

\begin{figure}[H]
\lstinputlisting[style=consola]{Img/consola-api.txt}
\caption{Salida de consola de las pruebas de la API.}
\label{fig:consola-api}
\end{figure}

\subsubsection{Integración en el navegador}
Se abrió \texttt{http://localhost:4200} en Google Chrome 157 y se jugó con un
script que pulsa las flechas y los botones de la interfaz.
\begin{longtable}{p{1cm}p{3.6cm}p{3.3cm}p{3.2cm}p{1.6cm}}
\caption{Integración en el navegador}\\
\toprule ID & Objetivo y procedimiento & Esperado & Obtenido & Estado\\\midrule\endfirsthead
\toprule ID & Objetivo y procedimiento & Esperado & Obtenido & Estado\\\midrule\endhead
I01 & Recorrer el modo manual con 54 flechas hasta la meta. & «Felicitaciones, buen trabajo.». & Meta alcanzada con ese mensaje. & PASS\\
I02 & Pulsar otra flecha en la meta. & El caminante no se mueve. & Posición sin cambios. & PASS\\
I03 & Resolver con la mano derecha. & Llegada a la meta con aviso. & «Has llegado al final». & PASS\\
I04 & Resolver con la mano izquierda. & Llegada a la meta con aviso. & «Has llegado al final». & PASS\\
I05 & Abrir la página a 390 px de ancho. & Sin desplazamiento horizontal. & 0 px de desborde. & PASS\\
I06 & Registrar errores de la página. & Ningún error del juego. & Solo un 404 de \texttt{favicon.ico}, que no afecta al juego. & PASS\\
\bottomrule
\end{longtable}

\subsubsection{Validación visual en el navegador}
\begin{figure}[H]
\centering
\begin{subfigure}[b]{0.72\linewidth}\centering
\captura[height=6cm]{go-menu.png}
\caption{Escritorio.}
\end{subfigure}\hfill
\begin{subfigure}[b]{0.26\linewidth}\centering
\captura[height=6cm,trim=0 400 0 0,clip]{go-celular-menu.png}
\caption{Celular (390 px).}
\end{subfigure}
\caption{Menú principal con los tres modos de juego.}
\label{fig:captura-menu}
\end{figure}

\begin{figure}[H]
\centering
\begin{subfigure}[t]{0.32\linewidth}\centering
\captura[width=0.99\linewidth,trim=150 0 150 60,clip]{go-manual-inicio.png}
\caption{Inicio.}
\end{subfigure}\hfill
\begin{subfigure}[t]{0.32\linewidth}\centering
\captura[width=0.99\linewidth,trim=150 0 150 60,clip]{go-manual-recorrido.png}
\caption{Tras 25 movimientos.}
\end{subfigure}\hfill
\begin{subfigure}[t]{0.32\linewidth}\centering
\captura[width=0.99\linewidth,trim=150 0 150 60,clip]{go-manual-meta.png}
\caption{Meta: «Felicitaciones, buen trabajo.».}
\end{subfigure}
\caption{Modo manual con el teclado.}
\label{fig:captura-manual}
\end{figure}

\begin{figure}[H]
\centering
\begin{subfigure}[t]{0.32\linewidth}\centering
\captura[width=0.99\linewidth,trim=150 0 150 60,clip]{go-right-inicio.png}
\caption{Antes de resolver.}
\end{subfigure}\hfill
\begin{subfigure}[t]{0.32\linewidth}\centering
\captura[width=0.99\linewidth,trim=150 0 150 60,clip]{go-right-recorrido.png}
\caption{En el recorrido.}
\end{subfigure}\hfill
\begin{subfigure}[t]{0.32\linewidth}\centering
\captura[width=0.99\linewidth,trim=150 0 150 60,clip]{go-right-meta.png}
\caption{Meta: «Has llegado al final».}
\end{subfigure}
\caption{Resolución automática con la mano derecha.}
\label{fig:captura-derecha}
\end{figure}

\begin{figure}[H]
\centering
\begin{subfigure}[t]{0.32\linewidth}\centering
\captura[width=0.99\linewidth,trim=150 0 150 60,clip]{go-left-inicio.png}
\caption{Antes de resolver.}
\end{subfigure}\hfill
\begin{subfigure}[t]{0.32\linewidth}\centering
\captura[width=0.99\linewidth,trim=150 0 150 60,clip]{go-left-recorrido.png}
\caption{En el recorrido.}
\end{subfigure}\hfill
\begin{subfigure}[t]{0.32\linewidth}\centering
\captura[width=0.99\linewidth,trim=150 0 150 60,clip]{go-left-meta.png}
\caption{Meta: «Has llegado al final».}
\end{subfigure}
\caption{Resolución automática con la mano izquierda.}
\label{fig:captura-izquierda}
\end{figure}

\begin{figure}[H]
\centering
\captura[height=8cm]{go-celular-juego.png}
\caption{Mano izquierda resuelta en un celular de 390 px.}
\label{fig:captura-celular}
\end{figure}

La voz se verificó con el manejador del backend en las pruebas; no se usó un
micrófono real.

% =====================================================================
\section{Análisis de Resultados: Métricas de Software}
% =====================================================================

\subsection{Líneas de código}
Se contabilizaron con el mismo criterio de los informes anteriores: líneas de
código (LOC), de comentario y en blanco.

\begin{table}[H]
\centering\small
\caption{Tamaño por archivo}
\begin{tabular}{lrrrr}
\toprule Archivo & LOC & Coment. & Blanco & Total\\\midrule
\texttt{backend/npc.go} & 463 & 0 & 59 & 522\\
\texttt{backend/walker.go} & 114 & 0 & 25 & 139\\
\texttt{backend/main.go} & 133 & 0 & 4 & 137\\
\texttt{backend/matrix.go} & 22 & 0 & 1 & 23\\
\textbf{Backend: producción} & \textbf{732} & \textbf{0} & \textbf{89} & \textbf{821}\\
\texttt{backend/parity\_test.go} & 128 & 0 & 3 & 131\\
\texttt{frontend/src/main.ts} & 65 & 0 & 0 & 65\\
\texttt{frontend/src/styles.css} & 13 & 1 & 12 & 26\\
\texttt{frontend/src/index.html} & 1 & 0 & 0 & 1\\
\textbf{Total} & \textbf{939} & \textbf{1} & \textbf{104} & \textbf{1\,044}\\
\bottomrule
\end{tabular}
\end{table}
El componente Angular y la hoja de estilos tienen pocas líneas porque su contenido
está en líneas largas (la hoja de estilos ocupa 15\,452 bytes).

\subsection{Otras métricas}
\begin{table}[H]
\centering\small
\caption{Métricas funcionales}
\begin{tabular}{p{7.8cm}r}
\toprule Métrica & Valor\\\midrule
Métodos de \texttt{Walker} / \texttt{NonPlayerCharacter} & 11 / 19\\
Pruebas de Go aprobadas & 3 / 3\\
Decisiones comparadas con el C\# & 320\\
Ticks hasta la meta: mano izquierda / derecha & 66 / 161\\
Latencia de un tick por la API (mediana / P95) & 3,45 / 4,94 ms\\
Paquete inicial de Angular (transferencia estimada) & 168,34 kB (50,95 kB)\\
Dependencias externas del backend & 0\\
\bottomrule
\end{tabular}
\end{table}
La latencia se midió en el mismo equipo, a través del proxy de desarrollo de
Angular; en un despliegue real dependerá de la red.

\subsection{Comportamientos observados}\label{sec:hallazgos}
\begin{enumerate}
  \item \textbf{Varios movimientos por tick en la mano derecha.} Por los \texttt{if}
        independientes, 15 de los 161 ticks mueven al caminante más de una casilla.
  \item \textbf{La voz no se bloquea en la meta.} El teclado deja de mover al
        caminante al llegar, pero la voz no, como en el original.
  \item \textbf{Mensaje de la mano derecha.} Esta versión muestra «Has llegado al
        final» en las dos manos, mientras que la versión JavaScript muestra
        «Haz llegado al final» en la mano derecha. Conviene confirmar con el
        código C\# cuál es el texto original.
  \item \textbf{Métodos sin uso.} \texttt{InitializeRapidMovement},
        \texttt{InitializePausedMovement} y \texttt{PausedMovementWithAI} se
        conservan aunque ninguna acción los usa.
  \item \textbf{Ritmo de los ticks.} El intervalo efectivo es de 50 ms más la
        latencia de cada petición, porque cada tick se pide al recibir el anterior.
\end{enumerate}

% =====================================================================
\section{Conclusiones y Recomendaciones}
% =====================================================================
\subsection{Conclusiones}
\begin{itemize}
 \item Se migró el Laberinto con IA a una arquitectura web con Go y Angular; el
       backend ejecuta toda la lógica con los mismos métodos de \texttt{Walker} y
       \texttt{NonPlayerCharacter}.
 \item Las 320 decisiones y los dos recorridos completos coinciden con el C\#
       original, tick por tick.
 \item La API valida sus entradas, separa las partidas por sesión y responde un
       tick en unos 3,5 ms, muy por debajo del intervalo de 50 ms.
 \item Se jugaron los tres modos en el navegador hasta la meta, en escritorio y en celular.
\end{itemize}

\subsection{Recomendaciones}
\begin{itemize}
 \item Unificar el mensaje de la mano derecha con el del original en ambas versiones.
 \item Agregar un \texttt{favicon.ico} para evitar el error 404 en la consola.
 \item Probar la voz con un micrófono real en varios navegadores.
 \item Liberar automáticamente las sesiones inactivas del backend.
\end{itemize}

% =====================================================================
\begin{thebibliography}{9}
\bibitem{ieee830} IEEE, \emph{IEEE Recommended Practice for Software Requirements Specifications}, IEEE Std 830-1998.
\bibitem{go} The Go Authors, \emph{The Go Programming Language Documentation}, \url{https://go.dev/doc/}.
\bibitem{angular} Google, \emph{Angular Documentation}, \url{https://angular.dev}.
\bibitem{mdn} Mozilla, \emph{MDN Web Docs: Web Speech API}, \url{https://developer.mozilla.org/docs/Web/API/Web_Speech_API}.
\bibitem{scrum} K. Schwaber y J. Sutherland, \emph{La Guía de Scrum}, 2020.
\bibitem{xp} K. Beck, \emph{Extreme Programming Explained: Embrace Change}, Addison-Wesley, 2004.
\end{thebibliography}

\end{document}
